\documentclass[11pt]{article}
\usepackage[margin=1in]{geometry}
\usepackage{amsmath,amssymb}
\usepackage{booktabs}
\usepackage{hyperref}
\usepackage{graphicx}
\usepackage{algorithm}
\usepackage{algpseudocode}
\usepackage{listings}

\title{lmscan: Statistical Feature Analysis for AI-Generated Text Detection and LLM Fingerprinting}
\author{Zacharie B \\
  \texttt{github.com/stef41/lmscan} \\
  Independent Researcher}
\date{April 2026}

\begin{document}
\maketitle

\begin{abstract}
We present \textsc{lmscan}, an open-source Python library for detecting AI-generated text and identifying the source large language model (LLM) using statistical feature analysis. Unlike transformer-based classifiers, \textsc{lmscan} computes 12 interpretable statistical features including burstiness, Zipf distribution deviation, hapax legomena ratio, and vocabulary richness. For model attribution, it scores text against 9 LLM family signatures based on characteristic vocabulary and structural patterns. The system operates in under 50ms with zero external dependencies, supports 5 languages, and provides full transparency into detection reasoning. We describe the feature extraction pipeline, the fingerprinting methodology, and discuss the tradeoffs between statistical and neural approaches to AI text detection.
\end{abstract}

\section{Introduction}

The widespread deployment of large language models (LLMs) has created an urgent need for AI text detection tools. Existing solutions---GPTZero, Originality.ai, Turnitin's AI module---are either closed-source, paid, or both. When these systems flag text as AI-generated, users cannot inspect why, leading to accountability concerns especially in educational settings \cite{liang2023gpt}.

We present an alternative approach: purely statistical feature analysis. Rather than training a classifier on labeled examples, \textsc{lmscan} computes a set of interpretable features that capture known distributional differences between human and machine-generated text. The key advantages are:

\begin{enumerate}
    \item \textbf{Transparency}: Every detection decision can be explained through individual feature contributions.
    \item \textbf{No training data}: The system does not require labeled examples to function (though calibration can improve domain-specific accuracy).
    \item \textbf{Zero infrastructure}: No GPU, API keys, or internet connection required.
    \item \textbf{Speed}: Full analysis in under 50 milliseconds.
\end{enumerate}

The obvious tradeoff is accuracy on adversarial text. A statistical approach will not match a fine-tuned transformer classifier on heavily paraphrased AI output. We argue, however, that for many practical use cases---screening student essays, auditing content pipelines, fingerprinting PR descriptions---the statistical approach offers a better balance of interpretability, cost, and reliability.

\section{Related Work}

Prior work on AI text detection falls into three categories:

\textbf{Watermarking}: Kirchenbauer et al.\ \cite{kirchenbauer2023watermark} propose embedding statistical signals during generation. This requires control over the generation process and does not work retroactively.

\textbf{Classifier-based}: DetectGPT \cite{mitchell2023detectgpt} uses perturbation-based likelihood comparisons. GPTZero and similar services train transformer classifiers on human/AI text pairs. These achieve high accuracy but are opaque and require substantial compute.

\textbf{Statistical}: Gehrmann et al.\ \cite{gehrmann2019gltr} analyze token-level probabilities using GLTR. Our work extends the statistical approach with a broader feature set and adds model attribution capabilities.

\section{Feature Extraction}

\textsc{lmscan} computes 12 features from input text. Each captures a different aspect of the distributional difference between human and LLM-generated text.

\subsection{Burstiness}

Human writing exhibits ``burstiness''---alternation between short, simple sentences and long, complex ones. LLM output tends toward unnaturally consistent sentence complexity. We measure burstiness as the coefficient of variation of per-sentence information entropy:

\begin{equation}
    B = \frac{\sigma(H_1, H_2, \ldots, H_n)}{\mu(H_1, H_2, \ldots, H_n)}
\end{equation}

where $H_i = -\sum_{w \in s_i} p(w) \log p(w)$ is the Shannon entropy of sentence $s_i$ computed over word frequencies within a sliding window.

\subsection{Zipf Distribution Deviation}

Natural language word frequencies follow Zipf's law: the $k$-th most common word appears with frequency proportional to $1/k^\alpha$, where $\alpha \approx 1$. LLM-generated text deviates from this distribution. We compute the Kolmogorov-Smirnov statistic between the observed rank-frequency distribution and the ideal Zipf curve:

\begin{equation}
    D_{KS} = \sup_x |F_{observed}(x) - F_{Zipf}(x)|
\end{equation}

\subsection{Hapax Legomena Ratio}

The proportion of words appearing exactly once (hapax legomena) in a text. Human writers use more unique, one-off words; LLMs tend to reuse a constrained vocabulary:

\begin{equation}
    R_{hapax} = \frac{|\{w : \text{count}(w) = 1\}|}{|V|}
\end{equation}

where $V$ is the vocabulary of the text.

\subsection{Vocabulary Richness (MATTR)}

Moving-Average Type-Token Ratio computed over sliding windows of 50 tokens. This normalizes for text length, which biases standard TTR:

\begin{equation}
    \text{MATTR} = \frac{1}{N-W+1} \sum_{i=1}^{N-W+1} \text{TTR}(w_i, \ldots, w_{i+W-1})
\end{equation}

\subsection{Additional Features}

\begin{itemize}
    \item \textbf{Sentence length variance}: Standard deviation of sentence lengths (words per sentence).
    \item \textbf{Punctuation entropy}: Shannon entropy of punctuation mark distribution.
    \item \textbf{N-gram repetition}: Frequency of repeated character-level and word-level trigrams.
    \item \textbf{Connective density}: Frequency of discourse connectives (``moreover'', ``furthermore'', ``however'').
    \item \textbf{Hedging ratio}: Frequency of hedging phrases (``I think'', ``perhaps'', ``it seems'').
    \item \textbf{Slop-word density}: Frequency of known LLM vocabulary markers (``delve'', ``tapestry'', ``paradigm'').
    \item \textbf{Paragraph structure}: Variance in paragraph length and regularity of paragraph breaks.
    \item \textbf{Overall entropy}: Character-level Shannon entropy of the full text.
\end{itemize}

\section{LLM Fingerprinting}

Beyond binary detection, \textsc{lmscan} identifies which LLM family likely generated a text. Each model has characteristic vocabulary and structural patterns:

\begin{table}[h]
\centering
\caption{Example vocabulary markers by model family}
\begin{tabular}{ll}
\toprule
\textbf{Model} & \textbf{Characteristic markers} \\
\midrule
GPT-4 & delve, tapestry, multifaceted, nuanced \\
Claude & I think, it's worth noting, I should mention \\
Llama & comprehensive, crucial, innovative, robust \\
Gemini & essentially, noteworthy, groundbreaking \\
Mistral & notably, fundamental, integral \\
\bottomrule
\end{tabular}
\end{table}

The fingerprinting module scores text against 9 model profiles (GPT-4, Claude, Gemini, Llama, Mistral, Qwen, DeepSeek, Cohere, Phi). Each profile consists of weighted vocabulary markers, hedging patterns, and structural signatures. The attribution confidence is computed as the normalized score differential between the best-matching and second-best-matching profiles.

\section{Implementation}

\textsc{lmscan} is implemented in pure Python with zero external dependencies. The architecture consists of:

\begin{itemize}
    \item \texttt{features.py}: 12 feature extractors
    \item \texttt{detector.py}: Weighted sigmoid classifier combining features into an AI probability score
    \item \texttt{fingerprint.py}: 9 model profile matchers
    \item \texttt{calibration.py}: Domain-specific threshold tuning via logistic regression on user-provided labeled data
    \item \texttt{languages.py}: Multilingual support (English, French, Spanish, German, Portuguese, CJK auto-detection)
\end{itemize}

The full analysis pipeline runs in under 50ms for typical-length texts (500--2000 words) on a single CPU core. The library is distributed via PyPI (\texttt{pip install lmscan}) and has 193 tests with CI on Python 3.9--3.13.

\section{Limitations and Future Work}

The statistical approach has clear limitations:

\begin{enumerate}
    \item \textbf{Adversarial robustness}: Paraphrasing or post-editing can defeat statistical detection.
    \item \textbf{Short texts}: Features are less reliable on texts under 100 words.
    \item \textbf{Domain sensitivity}: Thresholds tuned for academic text may not generalize to creative writing.
    \item \textbf{Model evolution}: As LLMs improve, their outputs become more human-like, reducing feature discriminability.
\end{enumerate}

Future directions include ensemble methods combining statistical and lightweight neural features, active learning for domain adaptation, and temporal analysis of how model fingerprints evolve across versions.

\section{Conclusion}

\textsc{lmscan} demonstrates that statistical feature analysis provides a viable, interpretable, and efficient approach to AI text detection and model attribution. While it does not match transformer-based classifiers on adversarial benchmarks, it offers unique advantages in transparency, cost, and deployment simplicity. The open-source release enables researchers and practitioners to inspect, extend, and calibrate the system for their specific needs.

\bibliographystyle{plain}
\begin{thebibliography}{9}

\bibitem{liang2023gpt}
Liang, W., et al.
\newblock GPT detectors are biased against non-native English writers.
\newblock \emph{Patterns}, 4(7), 2023.

\bibitem{kirchenbauer2023watermark}
Kirchenbauer, J., et al.
\newblock A watermark for large language models.
\newblock \emph{ICML}, 2023.

\bibitem{mitchell2023detectgpt}
Mitchell, E., et al.
\newblock DetectGPT: Zero-shot machine-generated text detection using probability curvature.
\newblock \emph{ICML}, 2023.

\bibitem{gehrmann2019gltr}
Gehrmann, S., et al.
\newblock GLTR: Statistical detection and visualization of generated text.
\newblock \emph{ACL Demo}, 2019.

\end{thebibliography}

\end{document}
